% GENERATED FILE. Edit canonical lessons, then run npm run generate:books.
% canonical-source-hash: fnv1a64:9bf0c8040451c664
% canonical-lessons: AR-C464-talawwuth, AR-C464-mulawwath, AR-C464-nufayat, AR-C464-lawwatha, AR-C464-faraza

\chapter{Pollution}
\label{ch:ar-pollution}
\hlchaptermodality{464}

\begin{chapteropening}
\textbf{By the end of this chapter:} \emph{I can say pollution, polluted, waste, to pollute and to sort.}

Complete the last lesson of chapter 464: I can say pollution, polluted, waste, to pollute and to sort.
\end{chapteropening}

\section[\texorpdfstring{\ar{تلوث} (talawwuth)}{تلوث (talawwuth)}]{\ar{تلوث} (talawwuth) — pollution}

\label{lesson:AR-C464-talawwuth}



\begin{quote}
Before the new one: say the Arabic for an emission, then the Arabic for carbon.
\end{quote}

\subsection*{You'll want to know: \ar{تلوث}}
\textbf{\ar{تلوث}} — \emph{talawwuth} — \textquotedblleft{}pollution\textquotedblright{}.

\begin{quote}
\ar{التلوث} \ar{في} \ar{المدينة} \ar{مشكلة} \ar{كبيرة}.
\end{quote}

\begin{quote}
\emph{at-talawwuth fī al-madīna mushkila kabīra.}
\end{quote}

\begin{quote}
Pollution in the city is a big problem.
\end{quote}

\begin{cousinweb}
\textbf{\ar{تلوث}} (\emph{talawwuth}) is built on the native Arabic root \emph{l-w-th}, \textquotedblleft{}to soil, to stain\textquotedblright{} (Lane). It is the verbal noun of Form V, \textquotedblleft{}becoming soiled\textquotedblright{}.
\end{cousinweb}

\begin{grammarlens}[title={What you've built}]
One more word for getting around.
\end{grammarlens}

\subsection*{Your turn}
\begin{itemize}
\raggedright
  \item \emph{Say it:} \emph{talawwuth}
  \item \emph{Say it:} \emph{talawwuth}, once more
  \item \emph{Say it:} \emph{at-talawwuth fī al-madīna mushkila kabīra.} — Pollution in the city is a big problem.
  \item \emph{Recall:} say the Arabic for an emission, then the Arabic for carbon, then say \emph{talawwuth} again
\end{itemize}

\begin{tcolorbox}[breakable,colback=teal!4,colframe=teal!35!black,title={Before you move on}]
What does \ar{تلوث} mean? (\textquotedblleft{}Pollution\textquotedblright{}.) Say it once more.
\end{tcolorbox}
\section[\texorpdfstring{\ar{ملوث} (mulawwath)}{ملوث (mulawwath)}]{\ar{ملوث} (mulawwath) — polluted}

\label{lesson:AR-C464-mulawwath}



\begin{quote}
Before the new one: say the Arabic for carbon, then the Arabic for pollution.
\end{quote}

\subsection*{You'll want to know: \ar{ملوث}}
\textbf{\ar{ملوث}} — \emph{mulawwath} — \textquotedblleft{}polluted\textquotedblright{}.

\begin{quote}
\ar{كان} \ar{ماء} \ar{النهر} \ar{ملوثا}.
\end{quote}

\begin{quote}
\emph{kāna māʾ an-nahr mulawwathan.}
\end{quote}

\begin{quote}
The river water was polluted.
\end{quote}

\begin{cousinweb}
\textbf{\ar{ملوث}} (\emph{mulawwath}) is built on the native Arabic root \emph{l-w-th}, \textquotedblleft{}to soil, to stain\textquotedblright{} (Lane). It is the Form II passive participle.
\end{cousinweb}

\begin{grammarlens}[title={What you've built}]
One more word for getting around.
\end{grammarlens}

\subsection*{Your turn}
\begin{itemize}
\raggedright
  \item \emph{Say it:} \emph{mulawwath}
  \item \emph{Say it:} \emph{mulawwath}, once more
  \item \emph{Say it:} \emph{kāna māʾ an-nahr mulawwathan.} — The river water was polluted.
  \item \emph{Recall:} say the Arabic for carbon, then the Arabic for pollution, then say \emph{mulawwath} again
\end{itemize}

\begin{tcolorbox}[breakable,colback=teal!4,colframe=teal!35!black,title={Before you move on}]
What does \ar{ملوث} mean? (\textquotedblleft{}Polluted\textquotedblright{}.) Say it once more.
\end{tcolorbox}
\section[\texorpdfstring{\ar{نفايات} (nufāyāt)}{نفايات (nufāyāt)}]{\ar{نفايات} (nufāyāt) — waste; refuse}

\label{lesson:AR-C464-nufayat}



\begin{quote}
Before the new one: say the Arabic for pollution, then the Arabic for polluted.
\end{quote}

\subsection*{You'll want to know: \ar{نفايات}}
\textbf{\ar{نفايات}} — \emph{nufāyāt} — \textquotedblleft{}waste; refuse\textquotedblright{}.

\begin{quote}
\ar{جمع} \ar{العمال} \ar{النفايات} \ar{من} \ar{الشارع}.
\end{quote}

\begin{quote}
\emph{jamaʿa al-ʿummāl an-nufāyāt min ash-shāriʿ.}
\end{quote}

\begin{quote}
The workers collected the waste from the street.
\end{quote}

\begin{cousinweb}
\textbf{\ar{نفايات}} (\emph{nufāyāt}) is built on the native Arabic root \emph{n-f-y}, \textquotedblleft{}to drive away, to cast out\textquotedblright{} (Lane). Waste is what is cast out.
\end{cousinweb}

\begin{grammarlens}[title={What you've built}]
One more word for getting around.
\end{grammarlens}

\subsection*{Your turn}
\begin{itemize}
\raggedright
  \item \emph{Say it:} \emph{nufāyāt}
  \item \emph{Say it:} \emph{nufāyāt}, once more
  \item \emph{Say it:} \emph{jamaʿa al-ʿummāl an-nufāyāt min ash-shāriʿ.} — The workers collected the waste from the street.
  \item \emph{Recall:} say the Arabic for pollution, then the Arabic for polluted, then say \emph{nufāyāt} again
\end{itemize}

\begin{tcolorbox}[breakable,colback=teal!4,colframe=teal!35!black,title={Before you move on}]
What does \ar{نفايات} mean? (\textquotedblleft{}Waste; refuse\textquotedblright{}.) Say it once more.
\end{tcolorbox}
\section[\texorpdfstring{\ar{لوث} (lawwatha)}{لوث (lawwatha)}]{\ar{لوث} (lawwatha) — to pollute}

\label{lesson:AR-C464-lawwatha}



\begin{quote}
Before the new one: say the Arabic for polluted, then the Arabic for waste.
\end{quote}

\subsection*{You'll want to know: \ar{لوث}}
\textbf{\ar{لوث}} — \emph{lawwatha} — \textquotedblleft{}to pollute\textquotedblright{}.

\begin{quote}
\ar{لوث} \ar{المصنع} \ar{ماء} \ar{البحيرة}.
\end{quote}

\begin{quote}
\emph{lawwatha al-maṣnaʿ māʾ al-buḥayra.}
\end{quote}

\begin{quote}
The factory polluted the water of the lake.
\end{quote}

Present: \ar{يلوث}, \emph{yulawwith}.

\begin{cousinweb}
\textbf{\ar{لوث}} (\emph{lawwatha}) is built on the native Arabic root \emph{l-w-th}, \textquotedblleft{}to soil, to stain\textquotedblright{} (Lane). It is Form II.
\end{cousinweb}

\begin{grammarlens}[title={What you've built}]
One more word for getting around.
\end{grammarlens}

\subsection*{Your turn}
\begin{itemize}
\raggedright
  \item \emph{Say it:} \emph{lawwatha}
  \item \emph{Say it:} \emph{lawwatha}, once more
  \item \emph{Say it:} \emph{lawwatha al-maṣnaʿ māʾ al-buḥayra.} — The factory polluted the water of the lake.
  \item \emph{Recall:} say the Arabic for polluted, then the Arabic for waste, then say \emph{lawwatha} again
\end{itemize}

\begin{tcolorbox}[breakable,colback=teal!4,colframe=teal!35!black,title={Before you move on}]
What does \ar{لوث} mean? (\textquotedblleft{}To pollute\textquotedblright{}.) Say it once more.
\end{tcolorbox}
\section[\texorpdfstring{\ar{فرز} (faraza)}{فرز (faraza)}]{\ar{فرز} (faraza) — to sort}

\label{lesson:AR-C464-faraza}



\begin{quote}
Before the new one: say the Arabic for waste, then the Arabic for to pollute.
\end{quote}

\subsection*{You'll want to know: \ar{فرز}}
\textbf{\ar{فرز}} — \emph{faraza} — \textquotedblleft{}to sort\textquotedblright{}.

\begin{quote}
\ar{فرزت} \ar{أمي} \ar{القمامة} \ar{في} \ar{المطبخ}.
\end{quote}

\begin{quote}
\emph{farazat ummī al-qumāma fī al-maṭbakh.}
\end{quote}

\begin{quote}
My mother sorted the rubbish in the kitchen.
\end{quote}

Present: \ar{يفرز}, \emph{yafriz}.

\begin{cousinweb}
\textbf{\ar{فرز}} (\emph{faraza}) is built on the native Arabic root \emph{f-r-z}, \textquotedblleft{}to separate, to set apart\textquotedblright{} (Lane). It is the plain verb.
\end{cousinweb}

\begin{grammarlens}[title={What you've built}]
One more word for getting around.
\end{grammarlens}

\subsection*{Your turn}
\begin{itemize}
\raggedright
  \item \emph{Say it:} \emph{faraza}
  \item \emph{Say it:} \emph{faraza}, once more
  \item \emph{Say it:} \emph{farazat ummī al-qumāma fī al-maṭbakh.} — My mother sorted the rubbish in the kitchen.
  \item \emph{Recall:} say the Arabic for waste, then the Arabic for to pollute, then say \emph{faraza} again
\end{itemize}

\begin{tcolorbox}[breakable,colback=teal!4,colframe=teal!35!black,title={Before you move on}]
What does \ar{فرز} mean? (\textquotedblleft{}To sort\textquotedblright{}.) Say it once more.
\end{tcolorbox}
